% Artículo VII, recuperación aditiva, revisión 04, 10-09-2026.
% RESULTADO_RECUPERADO: S15, 11c_gravedad_holonomia_rev6.tex:16--126.
% El cociclo canónico c27=(1,18) se conserva como antecedente del cómputo
% de rutas; este fragmento prueba su composición y representación.
% La memoria traslacional finita y la amplitud cosmológica s0 tienen
% tipos distintos. El cociclo no evalúa por sí solo una densidad local.

\section{Retorno de fase y representación de la memoria acumulada}
\label{vii:sec:retorno-memoria-gravedad}

La cronología de las rutas enriquecidas distingue el retorno de posición,
el retorno de orientación y el retorno de fase. La representación observable
identifica ciertos estados de llegada, mientras el registro acumulativo
conserva la secuencia recorrida. La representación geométrica debe
transportar ambos datos conjuntamente.

La conservación del complemento en
\eqref{g26:eq:archivo-completo-es} y la reconstrucción de residuos en
\eqref{g26:eq:accion-residuo-es} se aplican con este registro temporal
completo. Dos recorridos con igual fase observable pueden tener cociclos
de memoria distintos. Al transportar el lector radial se conservan
ambas coordenadas: el retorno periódico actúa en la fase, mientras la
concatenación incrementa el registro que figura en las fórmulas siguientes.

\subsection{Cociclo de inversión y cambio de hoja}
\label{vii:subsec:cociclo-memoria}

Sea \(F_{\rm obs}\) la representación cronológica del transporte TPK.
Sus bloques canónicos tienen la jerarquía
\[
\begin{array}{c|l}
27&\text{retorno posicional con inversión de orientación},\\
54&\text{retorno posicional y orientacional},\\
108&\text{retorno completo de la fase observable}.
\end{array}
\]
El registro cuenta las inversiones de orientación \(g(a)\) y los
cambios efectivos de hoja \(s(a)\). Para una ruta,
\begin{equation}
 c(\gamma)=\sum_{a\subset\gamma}(g(a),s(a)),\qquad
 c(\gamma_2\gamma_1)=c(\gamma_2)+c(\gamma_1).
 \label{vii:eq:cociclo-ruta}
\end{equation}
El cómputo canónico del bloque de longitud \(27\) expresa
\(c_{27}=(1,18)\). La composición conserva esos incrementos en
los cuatro bloques que completan el período observable.
En rutas orientadas hacia adelante el registro es acumulativo;
su extensión al grupoide de rutas utiliza \(\mathbb Z^2\) y
\(c(\bar\gamma)=-c(\gamma)\).

\begin{proposition}[Memoria de la vuelta observable completa]
\label{vii:prop:memoria-vuelta-completa}
Escribiendo \(\mathcal K_{27}=F_{\rm obs}^{27}\), la elevación
al registro, para \(x\in\mathcal O_{108}\) y \(z\in\mathbb Z^2\), satisface
\[
 \widetilde{\mathcal K}_{27}(x,z)
   =(\mathcal K_{27}x,z+(1,18)),\qquad
 \widetilde{\mathcal K}_{27}^{\,4}(x,z)
   =(x,z+(4,72)).
\]
Tras \(N\) vueltas completas, el incremento es \(N(4,72)\).
\end{proposition}

\begin{proof}
La proyección observable de \(\mathcal K_{27}\) tiene orden cuatro.
El cómputo canónico anterior fija \(c_{27}=(1,18)\) de manera constante
sobre esa órbita observable. La aditividad de
\eqref{vii:eq:cociclo-ruta} suma los cuatro incrementos y da
\(c_{108}=4(1,18)=(4,72)\).
La repetición aplica la misma ley de concatenación y produce
\(Nc_{108}\). El retorno de la primera coordenada conserva el
incremento de la segunda.
\end{proof}

\subsection{Representación afín del cociclo}
\label{vii:subsec:representacion-afin-cociclo}

La realización discreta utiliza un elemento \(R_4\) de orden cuatro en
\(\operatorname{Spin}^+(1,3)\), un vector temporal unitario \(u\)
fijado por su acción vectorial y una unidad de longitud \(\ell_0>0\).
La escala \(\ell_0\) conserva su procedencia en la sección dimensional.
Definimos
\begin{equation}
 \rho_{\rm C}^{\rm disc}(\gamma)=
 \bigl(R_4^{c_g(\gamma)},\,\ell_0c_s(\gamma)u\bigr)
 \in\operatorname{Spin}^+(1,3)\ltimes\mathbb R^{1,3}.
 \label{vii:eq:representacion-discreta-memoria}
\end{equation}
La potencia actúa sobre el vector mediante la representación vectorial
asociada del grupo de espín.

\begin{theorem}[Composición y amplitud traslacional del retorno]
\label{vii:thm:traslacion-memoria}
La aplicación \(\rho_{\rm C}^{\rm disc}\) conserva la identidad,
la concatenación y la inversión. En los retornos completos,
\begin{equation}
 \rho_{\rm C}^{\rm disc}(\gamma_{108})=(I,72\ell_0u),\qquad
 \rho_{\rm C}^{\rm disc}(\gamma_{108}^{\,N})=(I,72N\ell_0u).
 \label{vii:eq:traslacion-retorno-108}
\end{equation}
La inversión cambia el signo de la traslación.
\end{theorem}

\begin{proof}
El producto semidirecto es
\((R,a)(S,b)=(RS,a+Rb)\). Como \(R_4u=u\),
\[
\begin{split}
 \rho_{\rm C}^{\rm disc}(\gamma_2)
 \rho_{\rm C}^{\rm disc}(\gamma_1)
 &=
 \bigl(R_4^{c_g(\gamma_2)+c_g(\gamma_1)},
 \ell_0(c_s(\gamma_2)+c_s(\gamma_1))u\bigr)\\
 &=\rho_{\rm C}^{\rm disc}(\gamma_2\gamma_1).
\end{split}
\]
El cociclo de la identidad es cero y el de la inversa es el opuesto.
Para la vuelta completa se sustituyen \(c_{108}=(4,72)\) y
\(R_4^4=I\). La concatenación de \(N\) vueltas mantiene la parte
rotacional identidad y suma sus traslaciones.
\end{proof}

La amplitud \(72\ell_0\) es una longitud asociada a una ruta finita.
Una realización mediante cotetrada y conexión la transporta a su
holonomía afín. La contribución local a las ecuaciones gravitatorias
se obtiene después mediante la curvatura de esa conexión, la
variación de la acción material y la normalización de su densidad.
La representación conserva así el contenido acumulativo exacto que
debe intervenir en esa realización.
